% ════════════════════════════════════════════════════════════════════
%  thi hành — pha 1 sinh ra g_theta, cổng, rủi ro, tóm lại
% ════════════════════════════════════════════════════════════════════
\section{thi hành}

\begin{frame}[fragile]{$g_\theta$ ở đâu ra — pha 1}
  \framesubtitle{thi hành}
  \lead{cả timeline \texttt{c0}--\texttt{c7} giả định đã có $g_\theta$. nó được
        đào ở một pha riêng, chạy trước, trên một tập nhỏ hơn hai mươi lần.}

  \begin{columns}[t]
  \begin{column}{0.46\textwidth}
    \begin{bx}[title={pha 1 · $10{,}000$ truy vấn · llm đầy đủ}]
      \begin{dia}
a0  chẻ câu hỏi + bảng phủ       400
a1  vòng 0  ba mốc + 40 đoạn
            đơn lẻ            44.500
a2  vòng 1  greedy biên       27.000
a3  vòng 2  học chủ động     ~37.000
a4  vòng 3  xác nhận lại     ~15.000
                            ---------
                            ~123.900 token
      \end{dia}
      \vspace{0.25em}
      sản phẩm của pha 1 là \hl{$g_\theta$}, \no{không phải} nhãn cho tầng b.
    \end{bx}
  \end{column}
  \begin{column}{0.50\textwidth}
    \begin{bx}[title={pha 2 · $200{,}000$ truy vấn · đường rẻ}]
      \begin{dia}
c0  chẻ + bảng phủ               400
      thiếu mảnh -> THOÁT SỚM
c1  v(rỗng)                      500
c2  g_theta chấm 40 đoạn đơn       0
c3  v(R)                       3.500
c4  v(C)                      10.500
      -> chốt tau_q
c5  tìm theo tầng bằng g_theta     0
c6  đo v thật, kiểm tối tiểu  ~7.400
c7  ba nhánh
                              --------
                               ~22.300 token
      \end{dia}
      \vspace{0.25em}
      \dd{\textbf{rẻ hơn 5,5 lần.} sản phẩm của pha 2 mới là nhãn $(q \to T)$
          cho tầng b}
    \end{bx}
  \end{column}
  \end{columns}
\end{frame}

\begin{frame}{năm cổng — và chúng nằm trước 75\% chi phí}
  \framesubtitle{thi hành}
  \lead{không cam kết ngân sách trước. chia thành các chặng, mỗi chặng kết thúc
        bằng một \emph{cổng}: một phép đo rẻ quyết định có đi tiếp hay dừng.}

  \begin{center}
  \footnotesize
  \begin{tabular}[t]{@{}llrrl@{}}
    \toprule
    & \textbf{giai đoạn} & \textbf{chi phí} & \textbf{tích lũy} & \textbf{cổng chặn sau đó}\\
    \midrule
    \texttt{0}    & hạ tầng · index test có sẵn & \$0 & \$0 &\\
    \texttt{0.5}  & chạy thử bốn cổng trên test index & \$18 & \$18
      & \hl{cổng 0.5} — chặn \$1.100\\
    \texttt{1}    & hiệu chuẩn với chuẩn ngoài & \$7 & \$25
      & \hl{cổng 1} — $\mathrm{adq}$ có đáng tin?\\
    \texttt{2.0}  & dựng index cho phần pha 1 & \$74 & \$99 &\\
    \texttt{2}    & pha 1 đào nhãn & \$103 & \$202 &\\
    \texttt{3}    & huấn luyện $g_\theta$ & \$0 & \$202 & \hl{cổng 2} — vượt được sàn?\\
    \texttt{4}    & đo tỉ lệ $p$ trên phần giữ lại & \$0 & \$202 & \hl{cổng 3} — chặn \$918\\
    \texttt{5.0}  & dựng index cho phần pha 2 & \$815 & \$1.017 &\\
    \texttt{5}    & pha 2 $+$ ca chạy lại & \$103 & \$1.120 & \hl{cổng 4} $\times\,10$ lô\\
    \bottomrule
  \end{tabular}
  \end{center}

  \vspace{0.6em}
  \begin{keyb}
    tiêu trước \hl{cổng 0.5} là \$18; trước \hl{cổng 3} là \$202. thiết kế sai thì
    \yes{không bao giờ phải trả \$918 còn lại.}
    \vspace{0.2em}

    \dd{model sinh chạy trên máy mình nên \$0 tiền, nhưng ${\sim}53$ giờ gpu —
        ràng buộc thật là \hl{giờ gpu}, không phải tiền. $\$889$ trong bảng là
        dựng index, vì index dựng sẵn chỉ có cho tập \emph{test}}
  \end{keyb}
\end{frame}

\begin{frame}{những gì chưa chắc — và cổng nào đo nó}
  \framesubtitle{thi hành}
  \lead{năm con số dưới đây chưa ai đo. nói ra chúng, và nói luôn ai sẽ đo.}

  \begin{center}
  \begin{tabular}[t]{@{}>{\raggedright\arraybackslash}p{0.30\linewidth}
                        >{\raggedright\arraybackslash}p{0.21\linewidth}
                        >{\raggedright\arraybackslash}p{0.41\linewidth}@{}}
    \toprule
    \textbf{con số} & \textbf{trạng thái} & \textbf{ai đo}\\
    \midrule
    độ dài đoạn thật & giả định $250$ token
      & \texttt{0.b} — sai thì \no{mọi} bảng chi phí sai, và \no{không cổng nào bắt được}\\[0.35em]
    $p$ — tỉ lệ qua được đường rẻ & giả định $0{,}75$
      & \texttt{cổng 0.5} đo rẻ, \texttt{cổng 3} đo thật\\[0.35em]
    $45\%$ thoát sớm ở \texttt{c0} & suy từ musique
      & \texttt{2.1} — chắc chắn \emph{lạc quan}, vì ta cố ý dùng checkpoint chưa
        tinh chỉnh\\[0.3em]
    $\mathrm{adq}$ sai bao nhiêu & chưa biết
      & \texttt{cổng 1}, đối chiếu nhãn \texttt{supporting\_documents}\\[0.3em]
    cỡ rò rỉ nhãn vàng & chưa biết
      & \texttt{1.3} — chênh lệch giữa pool chỉ-bm25 và pool đầy đủ\\
    \bottomrule
  \end{tabular}
  \end{center}

  \vspace{0.5em}
  \begin{keyb}
    hai cổng \yes{miễn phí} phải chạy trước khi cam kết ngân sách:
    \hl{cổng $p$} ($500$ truy vấn giữ lại, 0 lượt llm) và
    \hl{chuẩn ngoài} ($500$ truy vấn test, \$7). chúng rẻ \emph{chính vì} nằm
    trước chi phí — bỏ qua để «chạy nhanh hơn» là bỏ đi toàn bộ giá trị của chúng.
  \end{keyb}
\end{frame}

\begin{frame}{tóm lại}
  \begin{columns}[t]
  \begin{column}{0.48\textwidth}
    \begin{keyb}[title={ba quyết định thiết kế}]
      \textbf{1}\; «đủ» có \emph{hai} thành phần. $v$ đo \emph{trả lời đúng},
      $\mathrm{adq}$ đo \emph{có bằng chứng}. chỉ dùng $v$ thì nhận cả đoạn 19.
      \vspace{0.4em}

      \textbf{2}\; tìm \emph{tập nhỏ nhất}, không xếp hạng điểm số. được cả
      \emph{họ} lời giải thích và một thước đo độ tin cậy; cái giá là mất
      đẳng thức kiểm tra.
      \vspace{0.4em}

      \textbf{3}\; chỉ $T$ đi qua ranh giới. hai tầng đo trên hai không gian
      khác nhau, nên điểm số không mang qua được.
    \end{keyb}
  \end{column}
  \begin{column}{0.48\textwidth}
    \begin{bx}[title={hai thứ còn phải chứng minh}]
      \textbf{$p$} — tỉ lệ truy vấn mà đường rẻ đi qua được. cả bảng chi phí treo
      vào nó, và nó \no{chưa được đo}.
      \vspace{0.4em}

      \textbf{$\mathrm{adq}$} — model nli loại đúng hay loại bừa. chưa biết, mà nó
      đang được giao quyền bỏ ${\sim}45\%$ truy vấn.
    \end{bx}
    \begin{bx}[title={bước tiếp theo}]
      chạy \texttt{0.a}--\texttt{0.j} trên test index (đã dựng sẵn, \$0),
      rồi \texttt{cổng 0.5} với \$18.
      \vspace{0.2em}

      \dd{github.com/pqvu0307/minimal\_set\_rag-graph\_learning}
    \end{bx}
  \end{column}
  \end{columns}
\end{frame}
